\chapter{Related Work\label{cha:related}}
% Drafted per POINT-SEQUENCE.md Ch.6 (points 3.1–3.7).

This chapter situates the thesis alongside existing systems and studies.
It covers the reference integration of RocksDB in stream processing, work on larger-than-memory window state, the tuning literature, and the methodology this thesis follows.

\section{RocksDB as a State Backend in Stream Processing\label{sec:rel-flink}}

% [3.1]
Apache Flink is a widely used stream processing system.
Flink ships RocksDB as its production out-of-core state backend~\cite{flink2021}.
This integration is the reference point for state that exceeds memory in stream processing.

% [3.2]
Flink's runtime differs from NebulaStream's in two ways that block direct transfer of its results.
Flink organizes its state access around checkpoints, which are periodic snapshots for fault tolerance.
Its data exchange is pull-based at the boundaries that Marcu and Bouvry analyze~\cite{marcu2022}.
NebulaStream pushes data through pipelines and shares one global buffer pool across all operators (Section~\ref{sec:bg-nes}).
Results from Flink therefore do not transfer directly to NebulaStream.

% [3.5]
A separate line of work studies RocksDB's parameters.
Jin et al.\ show that large-scale parameter tuning improves performance but leaves substantial write amplification even after optimization~\cite{jin2022}.
Ad\~ao et al.\ co-design the storage hierarchy below the \ac{KV} store and show that data placement across devices, not tuning alone, governs its I/O performance~\cite{adao2025}.
Wei et al.\ analyze the write amplification of RocksDB's compaction policies and show that even the best policy retains substantial write amplification~\cite{wei2025}.
Together, these studies show that configuration mitigates but does not remove the write and I/O costs.
This thesis therefore fixes a single research-informed configuration rather than searching over configurations (Section~\ref{sec:bg-rocksdb}).

\section{Larger-than-Memory Window State\label{sec:rel-windows}}

% [3.3]
WindFlow demonstrates larger-than-memory windowed processing in a single-node engine~\cite{frassinelli2025}.
Larger-than-memory window state is therefore addressed beyond Flink.

% [3.4]
AION keeps window state across memory and persistent storage so that late
events can still update windows that have already closed~\cite{esteves2020}.
It bounds storage growth with proactive caching and predictive cleanup.
Like this thesis, it pushes window state out of memory onto a storage device
(Section~\ref{sec:bg-windows}), though for late-event handling rather than a hard memory cap.

Neither WindFlow nor AION evaluates a push-based engine with one shared global buffer pool under hard memory caps.

\section{Evaluation Methodology\label{sec:rel-method}}

% [3.6]
Fragkoulis et al.\ survey the evolution of stream processing systems, including their state management~\cite{fragkoulis2024}.
The thesis pre-registers its hypotheses and reports negative results in full.
Chapter~\ref{cha:evaluation} applies this discipline throughout.

\section{Positioning\label{sec:rel-positioning}}

% [3.7]
To the best of the author's knowledge, no published work evaluates an out-of-core window-state backend within a push-based engine with a single shared global buffer pool under hard memory caps.
This thesis provides this evaluation.
Its contribution is evaluation-first; the backend exists to enable measurement.

The next chapter summarizes what the evaluation establishes and what remains open.
